% GENERATED FILE. Edit canonical lessons, then run npm run generate:books.
% canonical-source-hash: fnv1a64:35ce784667a9a0d8
% canonical-lessons: SA-C379-adhikah, SA-C379-paryaptah, SA-C379-pracurah, SA-C379-vipulah, SA-C379-vardhayati, SA-R379-review-372-379

\chapter{More Than Enough}
\label{ch:sa-more-than-enough}
\hlchaptermodality{379}

\begin{chapteropening}
\textbf{By the end of this chapter:} \emph{I can say how much will be enough if things go to plan.}

Complete the last lesson of chapter 379: I can say how much will be enough if things go to plan.
\end{chapteropening}

\section[\texorpdfstring{\sk{अधिकः} (adhikaḥ)}{अधिकः (adhikaḥ)}]{\sk{अधिकः} (adhikaḥ) — more; additional}

\label{lesson:SA-C379-adhikah}



\begin{quote}
Before the new one: say the Sanskrit for falling short, then the Sanskrit for dwindles.
\end{quote}

\subsection*{You'll want to know: \sk{अधिकः}}
\textbf{\sk{अधिकः}} — \emph{adhikaḥ} — \textquotedblleft{}more; additional\textquotedblright{}.

\begin{quote}
\sk{यदि} \sk{अधिकं} \sk{धनम्} \sk{अस्ति}, \sk{तर्हि} \sk{दरिद्राय} \sk{ददातु।}
\end{quote}

\begin{quote}
\emph{yadi adhikaṃ dhanam asti, tarhi daridrāya dadātu.}
\end{quote}

\begin{quote}
If there is extra money, give it to a poor man.
\end{quote}

\begin{cousinweb}
\textbf{\sk{अधिकः}} (\emph{adhikaḥ}) MW gives \emph{adhika} \textquotedblleft{}additional\textquotedblright{} and \textquotedblleft{}surpassing in number or quantity, more numerous\textquotedblright{}; MW gives no derivation here, so its origin is uncertain (MW).
\end{cousinweb}

\begin{grammarlens}[title={What you've built}]
One more word for saying what happens if.
\end{grammarlens}

\subsection*{Your turn}
\begin{itemize}
\raggedright
  \item \emph{Say it:} \emph{adhikaḥ}
  \item \emph{Say it:} \emph{adhikaḥ}, once more
  \item \emph{Say it:} \emph{yadi adhikaṃ dhanam asti, tarhi daridrāya dadātu.} — If there is extra money, give it to a poor man.
  \item \emph{Recall:} say the Sanskrit for falling short, then the Sanskrit for dwindles, then say \emph{adhikaḥ} again
\end{itemize}

\begin{tcolorbox}[breakable,colback=teal!4,colframe=teal!35!black,title={Before you move on}]
What does \sk{अधिकः} mean? (\textquotedblleft{}More; additional\textquotedblright{}.) Say it once more.
\end{tcolorbox}
\section[\texorpdfstring{\sk{पर्याप्तः} (paryāptaḥ)}{पर्याप्तः (paryāptaḥ)}]{\sk{पर्याप्तः} (paryāptaḥ) — sufficient}

\label{lesson:SA-C379-paryaptah}



\begin{quote}
Before the new one: say the Sanskrit for dwindles, then the Sanskrit for additional.
\end{quote}

\subsection*{You'll want to know: \sk{पर्याप्तः}}
\textbf{\sk{पर्याप्तः}} — \emph{paryāptaḥ} — \textquotedblleft{}sufficient\textquotedblright{}.

\begin{quote}
\sk{यदि} \sk{अन्नं} \sk{पर्याप्तम्}, \sk{तर्हि} \sk{अतिथिः} \sk{अपि} \sk{खादिष्यति।}
\end{quote}

\begin{quote}
\emph{yadi annaṃ paryāptam, tarhi atithiḥ api khādiṣyati.}
\end{quote}

\begin{quote}
If there is enough food, the guest will eat too.
\end{quote}

\begin{cousinweb}
\textbf{\sk{पर्याप्तः}} (\emph{paryāptaḥ}) is built on \emph{pari}, \textquotedblleft{}around\textquotedblright{}, and the root √āp, \textquotedblleft{}to reach, to obtain\textquotedblright{}; MW gives \emph{pary-āpta} \textquotedblleft{}sufficient for\textquotedblright{} (MW).
\end{cousinweb}

\begin{grammarlens}[title={What you've built}]
One more word for saying what happens if.
\end{grammarlens}

\subsection*{Your turn}
\begin{itemize}
\raggedright
  \item \emph{Say it:} \emph{paryāptaḥ}
  \item \emph{Say it:} \emph{paryāptaḥ}, once more
  \item \emph{Say it:} \emph{yadi annaṃ paryāptam, tarhi atithiḥ api khādiṣyati.} — If there is enough food, the guest will eat too.
  \item \emph{Recall:} say the Sanskrit for dwindles, then the Sanskrit for additional, then say \emph{paryāptaḥ} again
\end{itemize}

\begin{tcolorbox}[breakable,colback=teal!4,colframe=teal!35!black,title={Before you move on}]
What does \sk{पर्याप्तः} mean? (\textquotedblleft{}Sufficient\textquotedblright{}.) Say it once more.
\end{tcolorbox}
\section[\texorpdfstring{\sk{प्रचुरः} (pracuraḥ)}{प्रचुरः (pracuraḥ)}]{\sk{प्रचुरः} (pracuraḥ) — plentiful}

\label{lesson:SA-C379-pracurah}



\begin{quote}
Before the new one: say the Sanskrit for additional, then the Sanskrit for sufficient.
\end{quote}

\subsection*{You'll want to know: \sk{प्रचुरः}}
\textbf{\sk{प्रचुरः}} — \emph{pracuraḥ} — \textquotedblleft{}plentiful\textquotedblright{}.

\begin{quote}
\sk{यदि} \sk{वर्षा} \sk{प्रचुरा}, \sk{तर्हि} \sk{धान्यम्} \sk{अपि} \sk{प्रचुरं} \sk{भविष्यति।}
\end{quote}

\begin{quote}
\emph{yadi varṣā pracurā, tarhi dhānyam api pracuraṃ bhaviṣyati.}
\end{quote}

\begin{quote}
If the rain is plentiful, the grain will be plentiful too.
\end{quote}

\begin{cousinweb}
\textbf{\sk{प्रचुरः}} (\emph{pracuraḥ}) MW gives \emph{pracura} \textquotedblleft{}much, many, abundant, plentiful\textquotedblright{} and gives no derivation, so its origin is uncertain (MW).
\end{cousinweb}

\begin{grammarlens}[title={What you've built}]
One more word for saying what happens if.
\end{grammarlens}

\subsection*{Your turn}
\begin{itemize}
\raggedright
  \item \emph{Say it:} \emph{pracuraḥ}
  \item \emph{Say it:} \emph{pracuraḥ}, once more
  \item \emph{Say it:} \emph{yadi varṣā pracurā, tarhi dhānyam api pracuraṃ bhaviṣyati.} — If the rain is plentiful, the grain will be plentiful too.
  \item \emph{Recall:} say the Sanskrit for additional, then the Sanskrit for sufficient, then say \emph{pracuraḥ} again
\end{itemize}

\begin{tcolorbox}[breakable,colback=teal!4,colframe=teal!35!black,title={Before you move on}]
What does \sk{प्रचुरः} mean? (\textquotedblleft{}Plentiful\textquotedblright{}.) Say it once more.
\end{tcolorbox}
\section[\texorpdfstring{\sk{विपुलः} (vipulaḥ)}{विपुलः (vipulaḥ)}]{\sk{विपुलः} (vipulaḥ) — ample; abundant}

\label{lesson:SA-C379-vipulah}



\begin{quote}
Before the new one: say the Sanskrit for sufficient, then the Sanskrit for plentiful.
\end{quote}

\subsection*{You'll want to know: \sk{विपुलः}}
\textbf{\sk{विपुलः}} — \emph{vipulaḥ} — \textquotedblleft{}ample; abundant\textquotedblright{}.

\begin{quote}
\sk{यदि} \sk{धनं} \sk{विपुलम्}, \sk{तर्हि} \sk{ग्रामः} \sk{सुखी} \sk{भविष्यति।}
\end{quote}

\begin{quote}
\emph{yadi dhanaṃ vipulam, tarhi grāmaḥ sukhī bhaviṣyati.}
\end{quote}

\begin{quote}
If there is ample money, the village will be happy.
\end{quote}

\begin{cousinweb}
\textbf{\sk{विपुलः}} (\emph{vipulaḥ}) is built on \emph{vi} and \emph{pula}, which MW takes to be probably the same as \emph{pura}, \textquotedblleft{}full\textquotedblright{}; MW gives \emph{vi-pula} \textquotedblleft{}large, extensive, abundant\textquotedblright{} (MW).
\end{cousinweb}

\begin{grammarlens}[title={What you've built}]
One more word for saying what happens if.
\end{grammarlens}

\subsection*{Your turn}
\begin{itemize}
\raggedright
  \item \emph{Say it:} \emph{vipulaḥ}
  \item \emph{Say it:} \emph{vipulaḥ}, once more
  \item \emph{Say it:} \emph{yadi dhanaṃ vipulam, tarhi grāmaḥ sukhī bhaviṣyati.} — If there is ample money, the village will be happy.
  \item \emph{Recall:} say the Sanskrit for sufficient, then the Sanskrit for plentiful, then say \emph{vipulaḥ} again
\end{itemize}

\begin{tcolorbox}[breakable,colback=teal!4,colframe=teal!35!black,title={Before you move on}]
What does \sk{विपुलः} mean? (\textquotedblleft{}Ample; abundant\textquotedblright{}.) Say it once more.
\end{tcolorbox}
\section[\texorpdfstring{\sk{वर्धयति} (vardhayati)}{वर्धयति (vardhayati)}]{\sk{वर्धयति} (vardhayati) — increases (something)}

\label{lesson:SA-C379-vardhayati}



\begin{quote}
Before the new one: say the Sanskrit for plentiful, then the Sanskrit for ample.
\end{quote}

\subsection*{You'll want to know: \sk{वर्धयति}}
\textbf{\sk{वर्धयति}} — \emph{vardhayati} — \textquotedblleft{}increases (something)\textquotedblright{}.

\begin{quote}
\sk{यदि} \sk{वणिक्} \sk{मूल्यं} \sk{वर्धयति}, \sk{तर्हि} \sk{जनाः} \sk{अल्पं} \sk{क्रीणन्ति।}
\end{quote}

\begin{quote}
\emph{yadi vaṇik mūlyaṃ vardhayati, tarhi janāḥ alpaṃ krīṇanti.}
\end{quote}

\begin{quote}
If the merchant raises the price, people buy less.
\end{quote}

Past (he …): \sk{सः} \sk{मूल्यं} \sk{वर्धितवान्}, \emph{saḥ mūlyaṃ vardhitavān}.

\begin{cousinweb}
\textbf{\sk{वर्धयति}} (\emph{vardhayati}) comes from the root √vṛdh, \textquotedblleft{}to grow\textquotedblright{}, in its causative form; MW gives \emph{vardhita} \textquotedblleft{}increased, grown, augmented\textquotedblright{} (MW).
\end{cousinweb}

\begin{grammarlens}[title={What you've built}]
One more word for saying what happens if.
\end{grammarlens}

\subsection*{Your turn}
\begin{itemize}
\raggedright
  \item \emph{Say it:} \emph{vardhayati}
  \item \emph{Say it:} \emph{vardhayati}, once more
  \item \emph{Say it:} \emph{yadi vaṇik mūlyaṃ vardhayati, tarhi janāḥ alpaṃ krīṇanti.} — If the merchant raises the price, people buy less.
  \item \emph{Recall:} say the Sanskrit for plentiful, then the Sanskrit for ample, then say \emph{vardhayati} again
\end{itemize}

\begin{tcolorbox}[breakable,colback=teal!4,colframe=teal!35!black,title={Before you move on}]
What does \sk{वर्धयति} mean? (\textquotedblleft{}Increases (something)\textquotedblright{}.) Say it once more.
\end{tcolorbox}
\section[(dialogue)]{Looking back: from What Stands in the Way to More Than Enough}

\label{lesson:SA-R379-review-372-379}



\begin{quote}
Say the Sanskrit for an obstacle and a hindrance.
\end{quote}

\subsection*{Your turn}
Say these aloud:

\begin{itemize}
\raggedright
  \item an obstacle, a hindrance, hinders, holds up, removes
  \item permission, forbids, a restriction, an exception, compulsory
  \item attentive, carelessness, avoids, a mishap, ceases
  \item adversity, a calamity, destruction, a drought, remedies
  \item qualified, capable, controls, control, takes refuge in
  \item exceedingly, extremely, intense, completely, exceeds
  \item a tiny portion, in a small degree, a half, falling short, dwindles
  \item additional, sufficient, plentiful, ample, increases (something)
\end{itemize}

\begin{tcolorbox}[breakable,colback=teal!4,colframe=teal!35!black,title={Before you move on}]
An obstacle? (\textbf{\sk{विघ्नः}}, \emph{vighnaḥ}.) A hindrance? (\textbf{\sk{बाधा}}, \emph{bādhā}.) Hinders? (\textbf{\sk{बाधते}}, \emph{bādhate}.) Holds up? (\textbf{\sk{विलम्बयति}}, \emph{vilambayati}.) Removes? (\textbf{\sk{अपनयति}}, \emph{apanayati}.) Permission? (\textbf{\sk{अनुमतिः}}, \emph{anumatiḥ}.) Forbids? (\textbf{\sk{निषेधति}}, \emph{niṣedhati}.) A restriction? (\textbf{\sk{प्रतिबन्धः}}, \emph{pratibandhaḥ}.) An exception? (\textbf{\sk{अपवादः}}, \emph{apavādaḥ}.) Compulsory? (\textbf{\sk{अनिवार्यः}}, \emph{anivāryaḥ}.) Attentive? (\textbf{\sk{अवहितः}}, \emph{avahitaḥ}.) Carelessness? (\textbf{\sk{प्रमादः}}, \emph{pramādaḥ}.) Avoids? (\textbf{\sk{परिहरति}}, \emph{pariharati}.) A mishap? (\textbf{\sk{अनर्थः}}, \emph{anarthaḥ}.) Ceases? (\textbf{\sk{विरमति}}, \emph{viramati}.) Adversity? (\textbf{\sk{विपत्तिः}}, \emph{vipattiḥ}.) A calamity? (\textbf{\sk{आपद्}}, \emph{āpad}.) Destruction? (\textbf{\sk{विनाशः}}, \emph{vināśaḥ}.) A drought? (\textbf{\sk{अनावृष्टिः}}, \emph{anāvṛṣṭiḥ}.) Remedies? (\textbf{\sk{प्रतिकरोति}}, \emph{pratikaroti}.) Qualified? (\textbf{\sk{योग्यः}}, \emph{yogyaḥ}.) Capable? (\textbf{\sk{समर्थः}}, \emph{samarthaḥ}.) Controls? (\textbf{\sk{नियच्छति}}, \emph{niyacchati}.) Control? (\textbf{\sk{नियन्त्रणम्}}, \emph{niyantraṇam}.) Takes refuge in? (\textbf{\sk{आश्रयते}}, \emph{āśrayate}.) Exceedingly? (\textbf{\sk{अतीव}}, \emph{atīva}.) Extremely? (\textbf{\sk{अत्यन्तम्}}, \emph{atyantam}.) Intense? (\textbf{\sk{तीव्रः}}, \emph{tīvraḥ}.) Completely? (\textbf{\sk{नितराम्}}, \emph{nitarām}.) Exceeds? (\textbf{\sk{अतिक्रामति}}, \emph{atikrāmati}.) A tiny portion? (\textbf{\sk{लेशः}}, \emph{leśaḥ}.) In a small degree? (\textbf{\sk{मनाक्}}, \emph{manāk}.) A half? (\textbf{\sk{अर्धम्}}, \emph{ardham}.) Falling short? (\textbf{\sk{न्यूनः}}, \emph{nyūnaḥ}.) Dwindles? (\textbf{\sk{ह्रसति}}, \emph{hrasati}.) Additional? (\textbf{\sk{अधिकः}}, \emph{adhikaḥ}.) Sufficient? (\textbf{\sk{पर्याप्तः}}, \emph{paryāptaḥ}.) Plentiful? (\textbf{\sk{प्रचुरः}}, \emph{pracuraḥ}.) Ample? (\textbf{\sk{विपुलः}}, \emph{vipulaḥ}.) Increases (something)? (\textbf{\sk{वर्धयति}}, \emph{vardhayati}.)
\end{tcolorbox}
